\documentclass{article}
\usepackage{graphicx} % Required for inserting images
\usepackage{ulem} %for \sout -> strike through text

\title{BA Besprechungsnotizen}
\author{Bjarne Jehn}
\date{April 2026}

\begin{document}
	
\maketitle
%\section*{Besprechung am xx.xx.}
%\subsection*{Nächste Schritte}
%	\begin{itemize}
%		\item 
%	\end{itemize}
%	\subsection*{Fragen}
%	\begin{itemize}
%		\item 
%	\end{itemize}

\maketitle
\section*{Besprechung am 22.07.}
\subsection*{Nächste Schritte}
	\begin{itemize}
			\item Anmerkungen übernehmen
			\item Bennenung von minimizer länge m vs. minimizer min/M/..?
			\item frequency immer als funktion f(..) nutzen
			\item space consumption der Hashtabelle anpassen
		\end{itemize}
	\subsection*{Fragen}
	\begin{itemize}
			\item Hashtable/Lookup time? 
			\item warum werden binary constraints misachtet, wenn empirische Werte /100 geteilt werden? - numerische Fehler?
		\end{itemize}

\section*{Besprechung am 08.07.}
\subsection*{Nächste Schritte}
	\begin{itemize}
			\item  gurobi solver probieren - bisher nicht nötig
		\end{itemize}
	\subsection*{Fragen}
	\begin{itemize}
			\item rshash/test/hashtable.cpp korrigieren: Zeile 72 und 91 jeweils "value" statt "kmer\_value" und "value\_rev" statt "kmer\_value\_rev"
		\end{itemize}

\section*{Besprechung am 17.06.}
\subsection*{Nächste Schritte}
	\begin{itemize}
			\item \sout{im ilp unterscheiden zwischen variablen und konstanten} 
			\item \sout{erstmal nur ein level modellieren}
			\item \sout{ilp begründen schriftlich}
		\end{itemize}
	\subsection*{Fragen}
	\begin{itemize}
			\item 
		\end{itemize}

\section*{Besprechung am 10.06.}
\subsection*{Nächste Schritte}
	\begin{itemize}
			\item linearisiern
			\item oder einfach mal probieren
			\item für fixe werte mal durch einen solver laufen lassen und ergebnisse bewerten
			\item schonmal mit dem schreiben beginnen
			\item erst danach über code ausdenken
			\item git!
		\end{itemize}
	\subsection*{Fragen}
	\begin{itemize}
			\item Sollte man die Space Consumption nicht optimieren sondern eher möglichst genau modellieren, um dann eine obere Grenze wie z.B. 8Gb setzten zu können?
		\end{itemize}

\section*{Besprechung am 13.05.}
\subsection*{Nächste Schritte}
	\begin{itemize}
			\item git zugang
		\end{itemize}
	\subsection*{Fragen}
	\begin{itemize}
			\item Ist die Annahme, dass die Minimizer-Verteilung in den höheren Leveln ähnlich wie in Level 1 ist allgemein korrekt? - höchst wahrscheinlich schon.
			\item Wie kann ich festellen, wie viele von $4^m$ möglichen Minimizern vorkommen? Bzw. ist das überhaupt relevant?
			\item $p_M$ oder $p_K$ im ILP
			\item wie kann $p_K$ modeliert werden?
		\end{itemize}

\section*{Besprechung am 06.05.}
\subsection*{Nächste Schritte}
	\begin{itemize}
		\item gleichungen linearisieren, wenn möglich
		\item $p_M^2$ und $p_M^3$ modellieren
		\item space consumption modellieren
	\end{itemize}
\subsection*{Fragen}
	\begin{itemize}
		\item Ist der Zugriff auf R in den Unterschiedlichen Leveln relevant für die Laufzeit? - Denke schon, da dadurch höhere Level "bestraft" werden
		\item Ist die Zugriffszeit auf R bedeutend abhängig von m und t? - Wahrscheinlich nicht entscheidend aber könnte ebenfalls modelliert werden
	\end{itemize}

\section*{Besprechung am 29.04.}
\subsection*{Nächste Schritte}
	\begin{itemize}
		\item ILP für Laufzeit und Größe überlegen
		
		\item tex studio installieren - done
		
		\item erstmal nur lookup betrachten + (streaming locate? testen obs funktioniert)
		
		\item nach dem timen nochmal über alle positionen iterieren, damit der compiler nichts wegoptimiert
		
		\item überprüfen, ob tatsächlicher locate/lookup linear abhängig von minimizer verteilung ist
		
		\item in build.cpp gucken für implementierung
		\item streaming query lookup aus main - queries durch textstellen ersetzten, an denen minimizer mit frequenz i vorkommen
		
		\item überlegung: tatsächlich muss ja zunächst im Bitvektor und dann im Text nachgeschaut werden - zeiten variieren je nach system und die parameter können systemabhängig optimiert werden
	\end{itemize}

\subsection*{Fragen}
	\begin{itemize}
		\item BA Template? - kommt
		\item minimizer stats output - geklärt
		\item Warum kann ich \textit{.../rshash\$ .build/source/rshash --help} nicht ausführen aber \textit{.../rshash/build/source\$ ./rshash --help} schon? - fehlender backslash im PFad
		\item Unterscheidet sich mini\_stats auf dem dev branch von mini\_stats auf dem main branch? - nein
		\item was wäre eine sinnvolle Projektstruktur, wenn ich selber Code schreiben will? - es müsste möglich sein, programme aus rshash von BA directory aufzurufen
		Wäre es okay, wenn ich eine  Kopie des rshash repositories erstelle und diese in meinem git\_lab hochlade? - ja aber nicht nötig
		\item Fragen zu code:
		\begin{itemize}
			\item vieles passiert bereits in mini\_stats - ist das effizient gelöst? - ja
			\item falls ja, kann ich den Code nutzten und anpassen? - ja
			\item wäre es sinnvoll die minimizer komprimiert zu speichern? z.B. nur zwei bit pro Base oder gleich als Zahl? - jain, wird durch seqan bib bereits effizient gelöst
			\item muss ich für jedes m eine Datei mit den minimizern anlegen oder reicht der Arbeitsspeicher, um das ganze in einer Datenstruktur zu behalten? - Arbeitsspreicher sollte vorerst ausreichen. Für sehr große Genome bzw. Sammlungen von vielen großen Genomen wären dateien nötig. Eine Datei wäre vielleicht trotzdem sinnvoll, um später damit zu arbeiten.
		\end{itemize}
	\end{itemize}
\end{document}